\documentclass[11pt,a4paper]{article}
% ============================================================
% Shared preamble for Claude Experiments reports (Ukrainian, LuaLaTeX)
% Compile with: lualatex report_*_uk.tex
% Include this AFTER \documentclass in each report:
%   \documentclass[11pt,a4paper]{article}
%   % ============================================================
% Shared preamble for Claude Experiments reports (Ukrainian, LuaLaTeX)
% Compile with: lualatex report_*_uk.tex
% Include this AFTER \documentclass in each report:
%   \documentclass[11pt,a4paper]{article}
%   % ============================================================
% Shared preamble for Claude Experiments reports (Ukrainian, LuaLaTeX)
% Compile with: lualatex report_*_uk.tex
% Include this AFTER \documentclass in each report:
%   \documentclass[11pt,a4paper]{article}
%   \input{../preamble_uk}
% ============================================================
% Cyrillic via LuaLaTeX + fontspec (pdflatex has no T2A here).
\usepackage{fontspec}
\setmainfont{DejaVu Serif}
\setsansfont{DejaVu Sans}
\setmonofont{DejaVu Sans Mono}
\usepackage[a4paper,top=2.5cm,bottom=2.5cm,left=2.5cm,right=2.5cm]{geometry}
\usepackage{amsmath,amssymb}
\usepackage{graphicx}
\usepackage{subcaption}
\usepackage{booktabs}
\usepackage{xcolor}
\usepackage[colorlinks=true,allcolors=blue]{hyperref}
\usepackage{microtype}
\usepackage{float}
\usepackage{array}
\microtypesetup{expansion=false}

% Color scheme
\definecolor{colMicro}{HTML}{222222}
\definecolor{colExact}{HTML}{2878b5}
\definecolor{colPlacement}{HTML}{e07a1f}
\definecolor{colBoundary}{HTML}{429c74}
\definecolor{colCurrent}{HTML}{c23b55}
\definecolor{colMM}{HTML}{c23b55}

\providecommand{\SIG}{\textcolor{green!50!black}{\textbf{*}}}
\providecommand{\NS}{\textcolor{gray}{n.s.}}
\providecommand{\DRMS}{\Delta\mathrm{RMST}}
\providecommand{\ci}[2]{$[{#1},\;{#2}]$}

% ============================================================
% Cyrillic via LuaLaTeX + fontspec (pdflatex has no T2A here).
\usepackage{fontspec}
\setmainfont{DejaVu Serif}
\setsansfont{DejaVu Sans}
\setmonofont{DejaVu Sans Mono}
\usepackage[a4paper,top=2.5cm,bottom=2.5cm,left=2.5cm,right=2.5cm]{geometry}
\usepackage{amsmath,amssymb}
\usepackage{graphicx}
\usepackage{subcaption}
\usepackage{booktabs}
\usepackage{xcolor}
\usepackage[colorlinks=true,allcolors=blue]{hyperref}
\usepackage{microtype}
\usepackage{float}
\usepackage{array}
\microtypesetup{expansion=false}

% Color scheme
\definecolor{colMicro}{HTML}{222222}
\definecolor{colExact}{HTML}{2878b5}
\definecolor{colPlacement}{HTML}{e07a1f}
\definecolor{colBoundary}{HTML}{429c74}
\definecolor{colCurrent}{HTML}{c23b55}
\definecolor{colMM}{HTML}{c23b55}

\providecommand{\SIG}{\textcolor{green!50!black}{\textbf{*}}}
\providecommand{\NS}{\textcolor{gray}{n.s.}}
\providecommand{\DRMS}{\Delta\mathrm{RMST}}
\providecommand{\ci}[2]{$[{#1},\;{#2}]$}

% ============================================================
% Cyrillic via LuaLaTeX + fontspec (pdflatex has no T2A here).
\usepackage{fontspec}
\setmainfont{DejaVu Serif}
\setsansfont{DejaVu Sans}
\setmonofont{DejaVu Sans Mono}
\usepackage[a4paper,top=2.5cm,bottom=2.5cm,left=2.5cm,right=2.5cm]{geometry}
\usepackage{amsmath,amssymb}
\usepackage{graphicx}
\usepackage{subcaption}
\usepackage{booktabs}
\usepackage{xcolor}
\usepackage[colorlinks=true,allcolors=blue]{hyperref}
\usepackage{microtype}
\usepackage{float}
\usepackage{array}
\microtypesetup{expansion=false}

% Color scheme
\definecolor{colMicro}{HTML}{222222}
\definecolor{colExact}{HTML}{2878b5}
\definecolor{colPlacement}{HTML}{e07a1f}
\definecolor{colBoundary}{HTML}{429c74}
\definecolor{colCurrent}{HTML}{c23b55}
\definecolor{colMM}{HTML}{c23b55}

\providecommand{\SIG}{\textcolor{green!50!black}{\textbf{*}}}
\providecommand{\NS}{\textcolor{gray}{n.s.}}
\providecommand{\DRMS}{\Delta\mathrm{RMST}}
\providecommand{\ci}[2]{$[{#1},\;{#2}]$}


\title{\textbf{Довідник метрик: як порівнювати MicroMacro та Full Micro}\\[0.5em]
\large Які метрики коректні, які оманливі, і чому}
\author{Two-Scale CTMC Project}
\date{31 липня 2026}

\begin{document}
\maketitle

\begin{abstract}
Цей документ пояснює повну ієрархію метрик порівняння, що використовуються в
проєкті, демонструє математичний механізм кожної з них і визначає,
які метрики валідні за censoring, а які вносять систематичний bias.
Центральне повідомлення: \textbf{conditional mean та conditional median (лише
завершені прогони) є зміщеними щоразу, коли дві моделі мають різні completion
rates. RMST --- коректна первинна метрика.}
Усі приклади використовують реальні дані з каліброваного бенчмарку chain-of-cliques
та з експерименту абляції CE-01.
\end{abstract}

\tableofcontents
\clearpage

%% ============================================================
\section{Ядро проблеми порівняння}
%% ============================================================

Ми порівнюємо два симулятори на одній і тій самій стохастичній епідемічній моделі:
\begin{itemize}
  \item \textbf{Full Micro}: точний Gillespie CTMC на повній мережі.
        Ground truth.
  \item \textbf{MicroMacro (MM)}: двомасштабна апроксимація.
        Використовує агрегований macro hazard $\lambda_{ij} = \beta W_{ij}\,\bar{i}_i\,\bar{s}_j$
        для інтеркомунітарних передач.
\end{itemize}

Ключова величина --- $T_{\mathrm{last}}$: час, коли епідемія вперше
досягає останньої громади.  У кожній реплікації симуляція йде до
$T_{\mathrm{end}}$; якщо прибуття до того моменту не сталося, $T_{\mathrm{last}}$
є \emph{censored} (невідоме, відомо лише що воно $> T_{\mathrm{end}}$).

Фундаментальна складність: \textbf{дві моделі зазвичай мають різні
completion rates} (частка реплікацій, де $T_{\mathrm{last}} <
T_{\mathrm{end}}$), бо MM поширюється в середньому швидше за Micro через
Jensen-ефект.

%% ============================================================
\section{Аналіз відмінностей методів: дизайн перед метриками}
\label{sec:methodology}
%% ============================================================

Метрика має сенс лише всередині \emph{дизайну} порівняння.  Повторюваний
режим збою в цьому проєкті --- не вибір неправильної формули, а вимога до
метрики відповісти на питання, на яке вона не може відповісти, або читання числа,
що мовчки змішує ефект інтересу з confounder.  Цей розділ формулює
принципи, що мають керувати будь-яким порівнянням Micro-vs-MM.  Каталог метрик
(розділ~\ref{sec:catalogue}) тоді --- це просто набір інструментів, що
використовується \emph{всередині} цих принципів, а не їх замінник.

\subsection{Спершу вирішіть, яке питання ви ставите}

Три питання виглядають схоже, але вимагають різних дизайнів і дозволяють різні
висновки:

\begin{description}
  \item[\textbf{Q1 --- Існування.}] Чи взагалі MM відрізняється від Micro,
        поза Monte-Carlo шумом?  $\to$ значущість $\DRMS$ (bootstrap CI) або
        двовибірковий тест розподілів.  Потребує достатнього $N$; більше нічого.
  \item[\textbf{Q2 --- Магнітуда і знак.}] У якому напрямі й наскільки?
        $\to$ $\DRMS$ та його horizon-нормована форма, плюс
        completion gap.  Потребує censoring-aware оцінювачів і фіксованого горизонту.
  \item[\textbf{Q3 --- Механізм / атрибуція.}] \emph{Чому} виникає різниця,
        і яка властивість мережі її контролює?  $\to$ контрольована
        абляція, що перемикає одну апроксимацію за раз, плюс сурогатний
        параметричний тест.  Потребує ізоляції confounders, а \emph{не} лише
        великого ефекту.
\end{description}

Центральна методологічна помилка --- відповідати на \textbf{Q3 експериментом Q2}:
звітувати про великий, неконтрольований $\DRMS$ і \emph{наративно} приписувати
механізм (``premature leakage''), який дизайн ніколи не ізолював.  Магнітуда --- це
не механізм, а знак, що змінюється між сценаріями (як у $\beta$-ratio
sweep), є свідченням того, що число визначають confounders, а не названий механізм.

\subsection{Принцип ізоляції}

Щоб приписати різницю причині, кожна \emph{інша} причина має бути зафіксована.
Конкретно:
\begin{itemize}
  \item \textbf{Paired seeds.}  Запускати Micro і кожен варіант MM на тій самій
        реалізації мережі й тому самому RNG-потоці micro layer.  Парна
        різниця $\delta_r = T_r^{MM}-T_r^{Micro}$ тоді усуває between-run
        дисперсію, яка інакше домінує в CI.  CE-01 робить це;
        chain-of-cliques sweep --- ні.
  \item \textbf{По одному фактору за раз --- для атрибуції.}  Щоб приписати
        похибку \emph{механізму}, використовуйте ablation ladder (розділ~\ref{sec:ladder}),
        що перемикає одну апроксимацію на щабель.  Щоб відобразити \emph{параметричну
        залежність}, варіюйте параметри на сітці, але \emph{записуйте кожну коваріату},
        щоб її ефект можна було зрегресувати; ніколи не читайте одиничний
        one-factor sweep як каузальний.
  \item \textbf{Сумірні умови.}  Два сценарії порівнянні лише якщо
        їх горизонт і completion regime збігаються
        (розділи~\ref{sec:cens},~\ref{sec:horizon}).  Інакше числа не
        лежать на одній осі.
\end{itemize}

\subsection{Передумова 1: звітуйте completion, потім використовуйте censoring-aware оцінювачі}
\label{sec:cens}

Звітуйте $p_{\mathrm{Micro}}$, $p_{\mathrm{MM}}$ та їх gap \emph{перед} будь-яким
числом часу прибуття.  Щойно gap перевищує $\approx 5$--$10\%$, кожна
completed-runs-only метрика --- conditional mean, conditional median, $W_1$,
KS --- зазнає selection-bias і \emph{не} має нести headline
(розділ~\ref{sec:bias} показує, як це дає неправильний знак).  Використовуйте RMST /
survival curves, які включають цензуровані прогони на їх відомій нижній межі.  Коли
горизонти не можна зіставити, поверніться до Kaplan--Meier survival порівняння або
re-truncate обидві моделі до спільного completion level перед порівнянням.

\subsection{Передумова 2: сумірність горизонтів}
\label{sec:horizon}

RMST залежить від $T_{\mathrm{end}}$, тож сирий $\DRMS$ має сенс лише
\emph{всередині} фіксованого горизонту.  $\beta$-ratio sweep використовував
$T_{\mathrm{end}}\in[2000,20000]$; порівняння його сирих зсувів по $r$ порівнює
різні осі, і виживають лише \emph{знак} та horizon-нормовані величини.
Правило: всередині сценарію звітуйте сирий зсув; \emph{між} сценаріями
звітуйте $\DRMS/T_{\mathrm{end}}$, rel-$W_1$ або completion-probability gap
$\Delta p = p_{\mathrm{MM}}-p_{\mathrm{Micro}}$.

\subsection{Оберіть нормалізацію, незалежну від ефекту}
\label{sec:norm}

\mbox{rel-}$W_1 = W_1 / \mathbb{E}[T^{Micro}]$ ділить на величину, що лежить
на \emph{тій самій} часовій шкалі, що й чисельник, тож вона зменшує великі абсолютні
похибки ``безкоштовно'': довший ланцюг роздуває знаменник так само швидко, як і
чисельник, і $3.5\%$ можуть виглядати заспокійливо за побудовою.  Завжди
доповнюйте її шкалою, незалежною від середнього:
\[
  \frac{W_1}{\mathrm{IQR}(T^{Micro})}
  \qquad\text{або}\qquad
  \frac{W_1}{\Delta T_{\mathrm{hop}}},\quad
  \Delta T_{\mathrm{hop}} = \frac{\mathbb{E}[T_{\mathrm{last}}]}{L-1}.
\]
$W_1$, що становить $3.5\%$ середнього, але $40\%$ одного per-hop приросту,
не є малим у жодному операційно релевантному сенсі.

\subsection{Відокремте систематичний сигнал від стохастичного шуму}

Between-run (phase) шум домінує в неспарених порівняннях і в поточковій
$I_{\mathrm{tot}}(t)$ смузі.  Наслідки:
\begin{itemize}
  \item Надавайте перевагу \emph{парному} $\delta_i$ для Q1/Q3: він вимірює лише
        похибку апроксимації для спільної реалізації, усуваючи phase noise.
  \item Читайте IQR $I_{\mathrm{tot}}(t)$ як ``наскільки мінлива епідемія'',
        а не ``наскільки зміщений MM'' (розділ~\ref{sec:itot}).
  \item Широка смуга навколо двох майже збіжних медіан \emph{не} є свідченням
        великої похибки апроксимації.
\end{itemize}

\subsection{Атрибуція механізму: ablation ladder}
\label{sec:ladder}

Атрибуція потребує драбини, кожен щабель якої перемикає рівно одну апроксимацію, тож
приріст $\DRMS$ між сусідніми щаблями зараховується цій одній зміні:

\begin{center}
\begin{tabular}{llp{6.2cm}}
\toprule
Rung & Hazard $\times$ Placement & Що ізолює приріст \\
\midrule
Full Micro                 & exact / exact     & ground truth ($\DRMS=0$) \\
Oracle E/E                 & exact / exact-target & лише numerical/synchronisation \\
Oracle E/U                 & exact / uniform   & \textbf{import placement} error \\
Aggregate A/B              & aggregate / boundary & \textbf{rate abstraction} error \\
MicroMacro A/U             & aggregate / uniform & placement $+$ rate (deployed model) \\
\bottomrule
\end{tabular}
\end{center}

CE-01, проведений так, приписує основну частину похибки \emph{rate abstraction},
а placement значущий лише за точного hazard.  \textbf{Застереження:}
oracle-щаблі читають точний список ребер і стани вузлів; вони є \emph{діагностичними
верхніми межами}, а не розгортуваними моделями --- реальний MM може знати лише
агреговані параметри, ніколи точну структуру.

\subsection{Від атрибуції до розгортуваного параметричного сурогату}
\label{sec:surrogate}

Щойно драбина локалізує похибку до конкретної апроксимації, розгортуване
питання полягає в тому, чи може функція від \emph{відомих параметрів} ---
spectral gap $\lambda_2$, середній ступінь, boundary-degree статистики --- відтворити
oracle-корекцію \emph{без} точної структури.  Коректний тест --- не
``чи корелює похибка з $\lambda_2$ по мережах'' (ця кореляція
сконфаунджена всім, що ко-варіює з $\lambda_2$).  Він такий: \emph{чи
переміщує впорскування сурогата MM вгору по драбині до oracle-щабля, за всіх
інших факторів фіксованих?}  Щоб це підтримати, параметричний sweep має записувати
$\lambda_2$, ступінь, boundary статистики та горизонт на реалізацію, щоб
сурогат можна було підігнати й валідувати out-of-sample, а не підганяти до
тих самих corner cases, що його мотивували.

\subsection{Розкрийте калібрування}
\label{sec:calib}

Кожне поточне число bias припускає $\beta_{\mathrm{macro}}=\beta$ та
$W_{ij}=k$, встановлені \emph{by fiat}.  Частина $\DRMS$ тоді усувна
підбором $\lambda_{ij}$ (\emph{calibration} error), а частина невіддільна за
лише агрегованого стану (\emph{structural} error).  Звітуйте
oracle-калібрований зсув як structural floor і маркуйте кожне некаліброване
число як таке; інакше читач не може відрізнити властивість методу від
артефакту довільного вибору.

\subsection{Включіть negative control}
\label{sec:negctrl}

Достовірний аналіз різниці показує ефект \emph{зникаючим} там, де теорія
каже, що він має зникнути.  Топологія clique--bridge--leaf ($n=100$, $k=2$) є
confounder-free контролем; звітуйте її $W_1$ і KS \emph{з числами}, навіть коли
null.  Ефект, що кількісно визначений лише в corner cases (мале $k$, повільне
поширення, велике $\beta_{\mathrm{out}}$) і залишений якісним у чистому випадку,
ще не встановлений.

\subsection{Чекліст}

\begin{enumerate}
  \item Сформулюйте питання: Q1, Q2 чи Q3.  Узгодьте дизайн із ним.
  \item Спершу звітуйте completion; якщо gap $>10\%$, приберіть completed-only метрики
        з headline.
  \item Фіксуйте горизонт усередині порівняння; нормалізуйте (і розкривайте $T_{\mathrm{end}}$)
        між порівняннями.
  \item Додайте effect-independent нормалізацію поряд із rel-$W_1$.
  \item Використовуйте paired seeds і парний $\delta_i$, щоб зняти phase noise.
  \item Для тверджень про механізм показуйте ablation ladder, а не один великий зсув.
  \item Для тверджень про параметри фіксуйте коваріати або регресуйте їх, і
        валідуйте сурогат out-of-sample.
  \item Розкривайте калібрування; відокремлюйте structural від calibration error.
  \item Включайте negative control з числами.
\end{enumerate}

\begin{table}[H]
\centering
\renewcommand{\arraystretch}{1.3}
\begin{tabular}{p{6.2cm}p{7.4cm}}
\toprule
\textbf{Pitfall} & \textbf{Коректна практика} \\
\midrule
Conditional mean/median за нерівних completion & Звітуйте completion; використовуйте RMST / survival (розд.~\ref{sec:cens}) \\
Сирий $\DRMS$ порівняний за різних $T_{\mathrm{end}}$ & Нормалізуйте cross-scenario; сирий лише в межах горизонту (розд.~\ref{sec:horizon}) \\
rel-$W_1$ прочитаний як ``малий'' сам по собі & Додайте $W_1/\mathrm{IQR}$ або $W_1/\Delta T_{\mathrm{hop}}$ (розд.~\ref{sec:norm}) \\
Механізм наративований з одного великого зсуву & Ablation ladder, один toggle на щабель (розд.~\ref{sec:ladder}) \\
Похибка ``пояснена'' кореляцією з $\lambda_2$ & Фіксуйте коваріати / регресуйте; валідуйте сурогат (розд.~\ref{sec:surrogate}) \\
Числа bias з $\beta_{\mathrm{macro}}=\beta$ без маркування & Звітуйте structural floor; маркуйте uncalibrated (розд.~\ref{sec:calib}) \\
Чистий випадок залишено якісним & Negative control з числами $W_1$/KS (розд.~\ref{sec:negctrl}) \\
\bottomrule
\end{tabular}
\caption{Відображення поширених пасток аналізу на відповідний принцип.}
\label{tab:pitfalls}
\end{table}

%% ============================================================
\section{Каталог метрик}
\label{sec:catalogue}
%% ============================================================

\subsection{Completion Rate}

\[
  p_{\mathrm{complete}} = \frac{\#\{T_{\mathrm{last}} < T_{\mathrm{end}}\}}{N}
\]

\textbf{Що вимірює:} чи проходить епідемія весь ланцюг
у межах горизонту симуляції.  Залежить як від внутрішньої стохастичності
епідемії, так і від bias апроксимації MM.

\textbf{Приклад (clique baseline, 100 reps):}
\[
  p_{\mathrm{Micro}} = 0.68,\quad p_{\mathrm{MM}} = 0.82.
\]
MM завершує на 14\% більше прогонів --- прямий доказ, що MM поширюється швидше в середньому.

\textbf{Коректне використання:} Звітуйте першим.  Великий gap ($|p_{\mathrm{MM}} -
p_{\mathrm{Micro}}| > 0.1$) вже інформативний і означає, що всі інші
completed-run метрики будуть зміщеними (див.\ розділ~\ref{sec:bias}).

\textbf{Обмеження:} Це один бінарний поріг; він не квантифікує,
\emph{наскільки} швидше чи повільніше.

\subsection{Conditional Mean та Median (лише завершені прогони)}
\label{sec:bias}

\[
  \hat\mu^{\mathrm{cond}} = \frac{1}{|\mathcal{C}|}\sum_{r\in\mathcal{C}} T_r,
  \quad
  \mathcal{C} = \{r : T_r < T_{\mathrm{end}}\}.
\]

\textbf{Що вимірює:} середній (або медіанний) час прибуття
\emph{серед прогонів, що випадково завершилися}.

\paragraph{Чому воно зміщене, коли completion rates відрізняються.}
Припустимо, Jensen виконується: $F^{MM}(t) \geq F^{Micro}(t)$ для всіх $t$
(MM стохастично швидший, прибуває раніше).  Conditional median є
\[
  m^{Micro} = F^{-1}(0.5 \cdot p_{\mathrm{Micro}})
             = F^{-1}(0.34),
  \quad
  m^{MM}    = F^{-1}(0.5 \cdot p_{\mathrm{MM}})
             = F^{-1}(0.41).
\]
Ми порівнюємо $m^{MM}$ з $m^{Micro}$, але вони відповідають \emph{різним
перцентилям} безумовного розподілу (41-й проти 34-го).
Навіть коли $F^{MM}(t) \geq F^{Micro}(t)$, 41-й перцентиль Micro може
легко перевищити 41-й перцентиль MM, роблячи $m^{MM} > m^{Micro}$ ---
протилежне до того, що ми очікуємо.

\paragraph{Конкретний приклад.}
Clique baseline, 100 reps:
\[
  m_{\mathrm{cond}}^{Micro} = 1949.4,\quad
  m_{\mathrm{cond}}^{MM}    = 2020.2,\quad
  \text{shift} = +70.8 \quad \text{(MM здається пізнішим!)}.
\]
Але censoring-коректний RMST shift дорівнює $-75.6$ (MM раніший).
$+70.8$ --- це артефакт: додаткові 14\% completions MM --- це повільні прогони
поблизу $T_{\mathrm{end}} = 2500$, що тягнуть conditional median MM угору.

\textbf{Вирок: не використовуйте conditional mean чи median як первинну метрику,
коли completion rates відрізняються.}  Вони можуть показати неправильний знак.

\subsection{Restricted Mean Survival Time (RMST)}
\label{sec:rmst}

\[
  \mathrm{RMST}
  = \mathbb{E}\!\left[\min(T_{\mathrm{last}},\, T_{\mathrm{end}})\right]
  = \int_0^{T_{\mathrm{end}}} P(T_{\mathrm{last}} > t)\, dt.
\]

\textbf{Що вимірює:} очікуваний час до прибуття, обмежений зверху
$T_{\mathrm{end}}$.  Цензуровані прогони дають внесок $T_{\mathrm{end}}$ (їх нижню
межу), тож усі прогони включені без відкидання інформації.
Еквівалентно, RMST --- це площа під survival curve $S(t) =
P(T_{\mathrm{last}} > t)$ від $0$ до $T_{\mathrm{end}}$.

\textbf{Чому воно коректно обробляє censoring.}
Для цензурованого прогону ($T_r > T_{\mathrm{end}}$) ми знаємо, що справжній час
прибуття щонайменше $T_{\mathrm{end}}$.  Використання $\min(T_r, T_{\mathrm{end}}) =
T_{\mathrm{end}}$ як внеску --- це коректний likelihood-consistent
вибір для right-censored спостереження.

\textbf{Знаковий зсув:}
\[
  \DRMS = \mathrm{RMST}^{MM} - \mathrm{RMST}^{Micro}.
\]
\begin{itemize}
  \item $\DRMS < 0$: MM прибуває \emph{раніше} (передчасна пропагація).
  \item $\DRMS > 0$: MM прибуває \emph{пізніше} (надмірно консервативний).
  \item $\DRMS = 0$: немає систематичного bias на обраному горизонті.
\end{itemize}

\textbf{Зв'язок із Jensen.}
Нерівність Jensen передбачає $S^{Micro}(t) \geq S^{Macro}(t)$ для всіх $t$,
звідки $\DRMS \leq 0$: інтегрування нерівності survival дає
$\mathrm{RMST}^{MM} \leq \mathrm{RMST}^{Micro}$.  RMST тому є
коректною метрикою для \emph{перевірки} Jensen-прогнозу.

\paragraph{Як RMST обчислюється в цьому проєкті (код).}
З \texttt{extrascripts/benchmark\_censoring\_summary.py}:
\begin{verbatim}
restricted_micro = np.where(np.isfinite(micro), micro, horizon)
restricted_macro = np.where(np.isfinite(macro), macro, horizon)
rmst_shift = np.mean(restricted_macro - restricted_micro)
\end{verbatim}
Цензуровані прогони ($T_r = \mathrm{NaN}$ або $\infty$) замінюються на \texttt{horizon}
($T_{\mathrm{end}}$), потім зсув --- це середнє попарних різниць
по всіх $N$ реплікаціях.  Це еквівалентно різниці площ між
survival curves, оціненими емпіричним CDF.

\paragraph{Bootstrap довірчий інтервал.}
CI отримується ресемплінгом повної таблиці пар $(T_r^{Micro}, T_r^{MM})$
(включно з цензурованими рядками) із заміщенням, 1000 разів.  Пари
ресемплюються спільно (не незалежно), щоб зберегти within-run кореляції, якщо
дві моделі поділяють random seed.

\paragraph{Конкретні числа (калібрований clique benchmark, 100 reps):}
\begin{center}
\begin{tabular}{lrrrl}
\toprule
Scenario & $\mathrm{RMST}^{Micro}$ & $\mathrm{RMST}^{MM}$ & $\DRMS$ & 95\% CI \\
\midrule
baseline $k=5$  & 2089.3 & 2013.6 & $-75.6$  & $[-197.3,\; +34.8]$ \\
$k=10$          & 1326.0 & 1198.1 & $-127.9$ & $[-217.8,\; \mathbf{-39.7}]$ \\
$k=20$          & 879.5  & 738.7  & $-140.8$ & $[-190.6,\; \mathbf{-89.4}]$ \\
$L=15$          & 3342.1 & 3093.7 & $-248.4$ & $[-428.7,\; \mathbf{-51.2}]$ \\
sparse baseline & 2869.4 & 1640.5 & $-1228.9$& $[-1379.4,\; \mathbf{-1089.6}]$ \\
\bottomrule
\end{tabular}
\end{center}
Усі сценарії з достатнім completion показують $\DRMS < 0$ (жирне CI = значуще).
CI baseline $k=5$ включає нуль: не тому, що Jensen хибний, а тому, що
100 реплікацій недостатньо (SE $\approx 87$, зсув $\approx 0.87$ SE).

\textbf{Обмеження:} RMST залежить від $T_{\mathrm{end}}$.  Довший горизонт
робить RMST чутливішим до повільно-завершуваних прогонів.  Завжди звітуйте
горизонт поряд із RMST.

\subsection{Відстань Wasserstein-1 (W\texorpdfstring{$_1$}{1})}

\[
  W_1(P, Q) = \int_{-\infty}^{\infty} |F_P(t) - F_Q(t)|\, dt
            = \int_0^1 |F_P^{-1}(u) - F_Q^{-1}(u)|\, du.
\]

\textbf{Що вимірює:} середню абсолютну різницю між quantile
functions двох розподілів.  Еквівалентно, мінімальну очікувану відстань,
потрібну, щоб ``перевезти'' масу розподілу $P$ на $Q$.

\textbf{Relative W$_1$:}
\[
  \mathrm{rel}\text{-}W_1 = \frac{W_1}{\mathbb{E}[T^{Micro}]}.
\]
Це нормалізує на середнє Micro, даючи безрозмірну міру
розбіжності розподілів.

\textbf{Як воно використовується в цьому проєкті.}
$W_1$ обчислюється лише на \textbf{завершених} прогонах.  Тому воно зазнає
того самого censoring bias, що й conditional mean, коли completion rates
відрізняються.  Воно відповідає: ``Серед прогонів, які \emph{обидві} моделі
завершили, наскільки різними є розподіли часу прибуття?''  Це валідно для
порівняння \emph{форми розподілу}, але не вловлює
ефекту ``MM завершує більше прогонів''.

\textbf{Конкретні числа:}
\begin{center}
\begin{tabular}{lrrl}
\toprule
Scenario & $W_1$ & rel-$W_1$ & Note \\
\midrule
clique baseline $k=5$ & 67.0 & 3.5\% & малий, CI incl.\ zero \\
clique $k=10$         & 130.1 & 9.8\% & помірний \\
clique $k=20$         & 140.8 & 16.0\% & зростає з $k$ \\
sparse baseline       & ---  & 42.6\% & величезний; 99\% vs 100\% complete \\
\bottomrule
\end{tabular}
\end{center}

\textbf{Обмеження:} $W_1$ на завершених прогонах зміщений, коли completion rates
відрізняються.  Для sparse $k=1$ сценарію Micro завершує лише 2\%, а MM
завершує 19\%; $W_1$ не звітується, бо два підмножини не
порівнянні.

\textbf{Коли надавати перевагу $W_1$ над RMST:}
\begin{itemize}
  \item Коли обидві моделі мають 100\% completion (без censoring): $W_1$ і RMST
        дають узгоджену інформацію; $W_1$ чутливіший до відмінностей хвостів.
  \item Коли ви хочете порівняти \emph{форму} розподілу (не лише середнє):
        $W_1$ вловлює quantile-by-quantile відмінності, а не лише середнє.
  \item \textbf{Не} як первинну метрику, коли completion rates відрізняються.
\end{itemize}

\subsection{Статистика Kolmogorov--Smirnov}

\[
  D_{\mathrm{KS}} = \sup_t |F^{Micro}(t) - F^{MM}(t)|.
\]

\textbf{Що вимірює:} максимальну вертикальну відстань між двома CDF.
Асоційоване $p$-значення перевіряє, чи могли дві вибірки походити з
одного розподілу.

\textbf{Обмеження:} KS --- це \emph{якісний} тест (чи є різниця?),
а не \emph{кількісна} міра (наскільки велика різниця?).  Він також
обчислюється лише на завершених прогонах.  У цьому проєкті він використовується як
consistency check, а не первинна метрика.

\subsection{Per-Hop метрики часу прибуття: W\texorpdfstring{$_1$}{1}(T\texorpdfstring{$_i$}{i}) та парна похибка}

Усі метрики вище оперують на \emph{термінальному} часі прибуття $T_{\mathrm{last}}$.
Per-hop аналіз розкладає похибку шар за шаром.

\textbf{Per-hop Wasserstein $W_1(T_i)$:}
Обчисліть відстань $W_1$ між розподілом часів прибуття
$T_i^{Micro}$ та $T_i^{MM}$ для громади $i = 0,\ldots,L-1$
(по всіх реплікаціях, що досягли громади $i$).
Це відповідає: ``Де в ланцюзі пікує розподільна похибка?''

\textbf{Парна кумулятивна похибка:}
\[
  \delta_i = T_i^{MM} - T_i^{Micro} \quad \text{(та сама реплікація)}.
\]
Це обчислюється лише для прогонів, де обидві моделі досягли громади $i$
(парне спостереження).  Використовуючи той самий random seed, within-run кореляції
зберігаються, і оцінювач має значно нижчу дисперсію, ніж between-run
порівняння.

\textbf{Чому парна похибка --- найчутливіша метрика.}
Оскільки between-run мінливість домінує в невизначеності RMST,
парне порівняння усуває її повністю: кожен $\delta_i$ вимірює лише
\emph{похибку апроксимації} для того прогону, а не природну стохастичну
мінливість.  Ціна --- censoring: на останній громаді ($i=9$) з
68\% Micro completion існує лише 68 пар, а відсутні 32 --- це
\emph{найповільніші} прогони Micro --- саме ті з найбільшим $|\delta_i|$.

\textbf{Приклад (калібрований clique baseline):}
Медіанна парна похибка на шарі 7 становить $\approx -200$ (MM прибуває
на 200 одиниць часу раніше); на шарі 9 вона повертається до 0 через
censoring повільних прогонів Micro.  IQR становить $\pm 400$ --- велика, бо
кожен прогін зазнає різної локальної динаміки.

%% ============================================================
\section{Графік I\texorpdfstring{$_{\mathrm{tot}}$}{tot}(t): phase noise проти похибки апроксимації}
\label{sec:itot}
%% ============================================================

Загальна кількість інфікованих у момент $t$:
\[
  I_{\mathrm{tot}}(t) = \sum_{c=0}^{L-1} I_c(t).
\]
У постановці chain-of-cliques з $n=100$ вузлами на громаду, $I_{\mathrm{tot}}$
приблизно є східчастою функцією: кожна громада дає sigmoid-крок
висоти $n$.

\paragraph{Що насправді вимірює IQR-смуга на $I_{\mathrm{tot}}(t)$.}
У фіксований момент $t$ поточковий IQR по реплікаціях відображає:
\begin{enumerate}
  \item \textbf{Phase noise} ($\approx 90\%$ ширини): природну мінливість
        у тому, скільки громад інфіковано до моменту $t$.
        Різні прогони на різних ``сходинках'' драбини.
  \item \textbf{Похибку апроксимації} ($\approx 10\%$ ширини): систематичну
        різницю між Micro та MM у тому, наскільки швидко інфікуються громади.
\end{enumerate}

Оскільки phase noise домінує в ширині IQR, силует $I_{\mathrm{tot}}(t)$ ---
\textbf{не корисна міра якості апроксимації}.  Порівняння ширин
IQR-смуг Micro та MM говорить про природну мінливість епідемії, а не про
bias MM.

\paragraph{Коли графік $I_{\mathrm{tot}}(t)$ корисний.}
\begin{itemize}
  \item \textbf{Sanity check}: якщо дві медіани драматично різні
        (як у sparse baseline: $\Delta I \approx 400$ при $t=1500$),
        це прямий доказ великого bias.
  \item \textbf{Якісна ідентифікація режиму}: східчаста структура показує,
        скільки громад зазвичай інфіковано в кожен момент.
  \item \textbf{Не для}: квантифікації малих-до-помірних bias (clique case),
        чи порівняння форм розподілів.
\end{itemize}

%% ============================================================
\section{Ієрархія метрик та рекомендації}
%% ============================================================

\begin{table}[H]
\centering
\renewcommand{\arraystretch}{1.3}
\begin{tabular}{p{3.4cm}p{2.1cm}p{4.0cm}p{4.5cm}}
\toprule
\textbf{Metric} & \textbf{Обробляє censoring?} & \textbf{Використовувати коли} & \textbf{Не використовувати коли} \\
\midrule
Completion rate          & n/a        & Завжди звітувати першим             & --- \\
RMST / $\DRMS$           & \textbf{Так} & Первинна метрика; завжди валідна    & --- \\
$W_1(T_{\mathrm{last}})$ & Ні         & Обидві моделі $\approx100\%$ complete & Completion rates різняться $>5\%$ \\
Cond.\ mean shift        & Ні         & Completion rates $\approx$ рівні    & Completion rates різняться \\
Cond.\ median shift      & Ні         & Completion rates $\approx$ рівні    & Completion rates різняться \\
Paired $\delta_i$        & Частково    & Чутливе виявлення похибки per hop   & Останній шар (censoring повільних прогонів) \\
$W_1(T_i)$ per hop       & Ні         & Зростання похибки крізь ланцюг      & Перший/останній шари з малим completion \\
$I_{\mathrm{tot}}(t)$    & n/a        & Ідентифікація режиму, sanity check & Квантифікація похибки апроксимації \\
KS statistic             & Ні         & Лише якісний тест                   & Первинне кількісне порівняння \\
\bottomrule
\end{tabular}
\caption{Ієрархія метрик.  \textbf{Жирне}: рекомендована первинна метрика.}
\label{tab:hierarchy}
\end{table}

\paragraph{Рекомендований порядок звітування для кожного сценарію.}
\begin{enumerate}
  \item Completion rates $p_{\mathrm{Micro}}$, $p_{\mathrm{MM}}$ та їх
        різниця.
  \item $\DRMS$ з 95\% bootstrap CI.  Використовуйте це для перевірки Jensen і для
        оцінки практичного bias.
  \item rel-$W_1(T_{\mathrm{last}})$ \emph{лише якщо} обидва completion rates
        $\geq 0.9$.  Дає інформацію про форму розподілу поза середнім.
  \item Per-hop парна похибка $\delta_i$: медіана та IQR по реплікаціях.
        Для визначення, де в ланцюзі накопичується похибка.
  \item Медіана $I_{\mathrm{tot}}(t)$: лише для візуального контексту, не для
        кількісних висновків.
\end{enumerate}

%% ============================================================
\section{Розібраний приклад: абляція CE-01}
%% ============================================================

CE-01 використовує 100 \emph{парних} реплікацій (та сама мережа й seed для всіх
варіантів); $T_{\mathrm{end}} = 4000$; усі п'ять варіантів завершують 100\%
прогонів (без censoring).  Це ідеальна постановка: усі метрики валідні.

\paragraph{Таблиця первинних результатів.}
\begin{center}
\begin{tabular}{lrrrl}
\toprule
Variant & Completion & $\mathbb{E}[T_{\mathrm{last}}]$ & $\DRMS$ & 95\% CI \\
\midrule
Full Micro         & 1.00 & 2936.7 & $0$       & --- \\
Oracle E/E         & 1.00 & 2944.3 & $+7.5$    & $[-93.6,\;+120.3]$ \\
Oracle E/U         & 1.00 & 2775.5 & $-161.2$  & $[-266.8,\;-48.8]$ \\
Aggregate A/B      & 1.00 & 2735.5 & $-201.2$  & $[-302.4,\;-106.7]$ \\
MicroMacro (A/U)   & 1.00 & 2715.9 & $-220.8$  & $[-307.9,\;-120.2]$ \\
\bottomrule
\end{tabular}
\end{center}

Оскільки всі варіанти завершують 100\% прогонів і поділяють ті самі seeds,
$\DRMS = \mathbb{E}[T^{\mathrm{variant}}] - \mathbb{E}[T^{Micro}]$ точно
(без censoring, без between-run confounding).  Conditional mean та RMST
тут ідентичні.

\paragraph{Чому RMST усе одно має значення навіть за 100\% completion.}
Формулювання RMST чисто узагальнюється на цензурований випадок.
Використання RMST усюди --- навіть коли воно зводиться до звичайного середнього ---
забезпечує той самий код і ту саму інтерпретацію в усіх експериментах.
Це уникає введення окремих code paths, що могли б мовчки обчислити
неправильне, коли деякі прогони цензуровані.

%% ============================================================
\section{Статистична потужність та обсяг вибірки}
%% ============================================================

Поширене питання: скільки реплікацій потрібно, щоб виявити заданий
$\DRMS$ з 80\% потужністю на 5\% рівні значущості?

Для одновибіркового $t$-тесту на $\Delta_r = \min(T_r^{MM}, T_{\mathrm{end}}) -
\min(T_r^{Micro}, T_{\mathrm{end}})$ зі стандартним відхиленням $\sigma_\Delta$:
\[
  N \approx \left(\frac{(z_{0.025} + z_{0.20})\, \sigma_\Delta}{\delta}\right)^2
           = \left(\frac{2.80\, \sigma_\Delta}{\delta}\right)^2,
\]
де $\delta = |\DRMS|$ --- цільовий розмір ефекту.

\paragraph{Калібрований clique baseline ($k=5$).}
З bootstrap: $\sigma_\Delta \approx 590$ (std попарних різниць),
$|\DRMS| = 75.6$.  Потрібне $N \approx (2.80 \times 590 / 75.6)^2 \approx 480$.
Поточні 100 реплікацій дають потужність $\approx 25\%$ --- узгоджено з CI,
що включає нуль.

\paragraph{Калібрований clique $k=10$.}
$\sigma_\Delta \approx 450$, $|\DRMS| = 127.9$.
Потрібне $N \approx (2.80 \times 450 / 127.9)^2 \approx 97$.
100 реплікацій якраз достатньо; CI виключає нуль.

\paragraph{CE-01 (paired).}
Спарювання реплікацій (той самий seed) драматично зменшує $\sigma_\Delta$
(between-run дисперсія усунена).  100 парних реплікацій достатньо
для всіх чотирьох контрастів із $|\DRMS| \geq 160$.

\textbf{Практичне правило:} для chain benchmark без спарювання
використовуйте щонайменше $N = 300$ реплікацій, коли $|\DRMS| < 100$.

%% ============================================================
\section{Підсумок}
%% ============================================================

\begin{enumerate}
  \item \textbf{RMST та $\DRMS$} --- коректні первинні метрики.
        Вони обробляють censoring, прямо пов'язані з нерівністю Jensen,
        і мають добре зрозумілі статистичні властивості.

  \item \textbf{Conditional mean/median} (лише завершені прогони) вносить
        selection bias щоразу, коли completion rates відрізняються.  Воно може показати
        \emph{неправильний знак} (напр., clique baseline: conditional median $= +70.8$,
        тоді як RMST $= -75.6$).

  \item \textbf{Відстань W$_1$} валідна для порівняння форми розподілу,
        коли обидві моделі мають майже-100\% completion.  Використовуйте її поряд із RMST,
        а не замість нього.

  \item \textbf{Поточкові графіки $I_{\mathrm{tot}}(t)$} відображають phase noise,
        а не похибку апроксимації.  Корисні для візуального контексту (особливо
        sparse case), але не для квантифікації bias.

  \item \textbf{Per-hop парна похибка $\delta_i$} --- найчутливіший
        діагностичний засіб: вона усуває between-run дисперсію, використовуючи paired seeds.
        Інтерпретуйте значення останнього шару обережно через censoring
        повільних прогонів Micro.

  \item \textbf{100 реплікацій} недостатньо для clique baseline
        ($k=5$), щоб виявити Jensen bias на 5\% значущості.
        Потрібно $\approx 400$--$500$ реплікацій або paired seeds.
\end{enumerate}

\end{document}
